% AUTORIDAD DE FUENTE:
% expose_VIII_body.tex, líneas 2031--2063; índices de escaneo 302--304;
% folios de fuente 291--293. La corrección material de la capa francesa
% se conserva reversiblemente mediante \sisi.

\medskip
\noindent 6.3. En lo que sigue suponemos que
$G=\underline G_{mS}$. Toda extensión $\overline G$ de
$\mathbb Z_{S,T}$ por $G=\underline G_{mS}$ es entonces representable
(6.2.4) y es, naturalmente, lisa sobre $S$, por ser una extensión del
esquema en grupos étale $\mathbb Z_{S,T}$ por el esquema en grupos liso
$\underline G_{mS}$. Además, dar un torsor $F$ bajo $\underline G_m$
equivale a dar un haz invertible $\mathbb L$ sobre $S$ (asociando a
este último el torsor de las secciones invertibles de $\mathbb L$).
Así, $\overline G$ queda descrito por un par $(\mathbb L,s)$, donde
$\mathbb L$ es un haz invertible sobre $S$ y $s$ es una trivialización
de $\mathbb L$ sobre $U=S\setminus T$, es decir, una sección invertible
de $\mathbb L|U$.

Cuando $U$ es esquemáticamente denso en $S$ y $S$ es localmente
noetheriano, o bien es reducido y su conjunto de componentes
irreducibles es localmente finito, dar tal par equivale a dar un
\underline{divisor $D$ sobre $S$ tal que}
\begin{equation}\tag{6.3.1}
 \operatorname{supp}D\subset T,
\end{equation}
asociando a dicho divisor el par
\begin{equation}\tag{6.3.2}
 \bigl(\underline O_S(-D),\sisi{s_D}{s_{-D}}|U\bigr)
\end{equation}
formado por el haz invertible $\underline O_S(-D)$ asociado a $-D$
(EGA IV 21.2.8.1) y la restricción a $U$ de su sección racional
canónica $s_{-D}$ (EGA IV 21.2.11 y 20.2.11(ii)). En cualquier caso,
sin hipótesis alguna sobre $S$, un divisor $D$ sobre $S$ que satisface
(6.3.1) define mediante (6.3.2) un par $(\mathbb L,s)$ y, por tanto,
una extensión $\overline G$ de $\mathbb Z_{S,T}$ por
$\underline G_m$. La llamaremos el \underline{modelo de
Néron--Raynaud de $\underline G_m$} asociado al divisor $D$ y al
subesquema cerrado $T$ de $S$ que satisface (6.3.1). Si
$T=\operatorname{supp}D$, omitiremos $T$ de la terminología.

%%%%%%%%%%  scan idx 303  |  folio 292  %%%%%%%%%%

Es inmediato que la formación de esta extensión $\overline G$ conmuta
con todo cambio de base $f:S'\to S$ para el cual esté definido el
divisor imagen inversa $f^*(D)$ (EGA IV 21.4.2), y en particular con
todo cambio de base \underline{plano}.

\medskip
\noindent 6.3.3. La descripción general de 6.2.1 de las secciones de
$\overline G$ muestra inmediatamente que pueden identificarse con las
secciones $f$ del haz $\mathcal M_S^*$ de secciones meromorfas regulares
sobre $S$ (EGA IV 20.1.8) cuyo divisor $\operatorname{div}(f)$ es,
localmente para la topología de Zariski, un múltiplo entero de $D$. Si
$V$ es un abierto de $S$ en el que la restricción de
$\operatorname{div}(f)$ coincide con la de $nD$, donde
$n\in\mathbb Z$, entonces la sección de $\mathbb Z_{S,T}$ definida por
$f$ es simplemente la imagen en (6.1.2) de la sección constante de
$\mathbb Z_S$ de valor $n$. La descripción precedente del grupo de
secciones de $\overline G$ sobre $S$ sigue siendo válida sobre todo
abierto de $S$ y, más generalmente, para el grupo de puntos de
$\overline G$ con valores en cualquier $S$-esquema plano $S'$. Como
$\overline G$ mismo es plano (e incluso liso) sobre $S$, esta
descripción vuelve a dar una caracterización completa del modelo de
Néron--Raynaud de $\underline G_m$ definido por $D$ y
$T\supset\operatorname{supp}D$.

\medskip
\noindent 6.4. Supongamos ahora que $S$ es \underline{normal}, que su
conjunto de componentes irreducibles es localmente finito (de modo que
$S$ es una suma de esquemas normales integrales), que $D$ es un divisor
positivo cuyas multiplicidades son todas iguales a $1$ (EGA IV 21.6),
que $T=\operatorname{supp}D$ y que $T$ es localmente irreducible, es
decir, que sus componentes irreducibles son disjuntas. Entonces, para
toda función racional $f$ sobre $S$ que sea regular sobre
$U=S\setminus T$, el divisor $\operatorname{div}(f)$ es localmente un
múltiplo de $D$. Por tanto, en este caso existe un isomorfismo canónico
de grupos

%%%%%%%%%%  scan idx 304  |  folio 293  %%%%%%%%%%

\begin{equation}\tag{6.4.1}
 \Gamma(S,\overline G)\xrightarrow{\ \sim\ }
 \Gamma(U,\underline G_{mU}).
\end{equation}
Supongamos además que $T$ es geométricamente unirrama y localmente
integral, hipótesis estable por cambio de base liso (EGA IV 17.5.7).
Como todas las hipótesis son entonces estables por cambio de base, de
(6.4.1) se deduce que \underline{para todo $S'$ liso sobre $S$ existe
un isomorfismo canónico de grupos, funtorial en $S'$,}
\begin{equation}\tag{6.4.2}
 \mathrm{Hom}_S(S',\overline G)\xrightarrow{\ \sim\ }
 \mathrm{Hom}_U(S'_U,\underline G_{mU})
 \simeq\Gamma(S'_U,\underline O_{S'_U}^*).
\end{equation}
Como $\overline G$ es él mismo liso sobre $S$, este isomorfismo
caracteriza el esquema en grupos $\overline G$ sobre $S$ salvo
isomorfismo único. También muestra que $\overline G$ desempeña el papel
de un ``modelo de Néron'' para el esquema en grupos
$\underline G_{mU}$ sobre el abierto $U$ de $S$. Junto con el hecho de
que fue introducido y utilizado por primera vez por M.~Raynaud, esto
justifica la terminología adoptada en 6.3 para $\overline G$.
